\section{Prisoner of the System}

I used to believe that novels such as \emph{Brave New World} or \emph{1984} were to serve as warnings to us, as representing something to be avoided. Unfortunately, it seems their intention is to prepare us for the shape of things to come. We can see the TV series \emph{The Prisoner in the same way.}

In the series, \#6 is a prisoner in a mock village, and \#2 --- whose identity usually changed from week to week is the visible leader. The goal of every \#2 is to extract a confession from \#6 and integrate him into the village. \#6 on the other hand, resists all such attempts. \#2 only implements policy; the identity of \#1 who actually creates policy is left unknown.

\#6 holds to the principle: ``I am a man, not a number.'' We can regard this as a metaphor for a transcendent principle, since his self-identity as a man transcends his identity as an anonymous individual in the village. In other words, he asserts his identity as a Person rather than an Individual, a concept with origins in Thomas Aquinas. For Evola,

\begin{quotex}
A person is an individual who is differentiated through his qualities, endowed with his own face, his \emph{proper nature}, and a series of attributes that make him who he is ... that make him fundamentally \emph{unequal}.
\end{quotex}

\#2 is a modernist, since he holds no transcendent principles.

Therefore, he can employ all means of social coercion --- drugs, seduction, opprobrium, threats, bribes --- designed to break down \#6. Because \#2 holds no principles, there is no possibility of dialogue between \#6 and \#2. Instead, their conversations are nothing but mutual posturing, so every episode devolves into a battle of wills.

Concerning the purpose of social institutions, Evola writes:

\begin{quotex}
The perfection of the human being is the end to which every healthy social institution must be subordinated. ... This perfection must be conceived on the basis of a process of individuation and progressive differentiation.
\end{quotex}


To the contrary, for the village, the end of the human being is to fit into the social structure without disruption. The villagers are under constant surveillance, a trend seen today from the control of finance and medicine sectors which will make those records available to authorities, to public CCTV and the street, and even GPS tracking of citizens. Ideas that would in the past have been restricted to a small circle of initiates are now publicly available on the Internet.

The illusion of dissent is also part of the village. When coercion and deception fail, \#2 participates in a ``democratic'' election, which, of course, will not result in any real change. Any and all attempts by \#6 to influence the villagers fail. Without any transcendent principles themselves, the villagers are motivated by fear and sentiment. Speeches in the public square, editorials, and personal conversations are all ineffectual. Ultimately, \#6 does find some allies and the behind-the-scenes superstructure of the village is destroyed. Revolution is the only solution.

Of course, that cannot happen today. First of all, because, while the \#2's are legion, there is no \#1 in total control, hence no single point of vulnerability. That is because there is no real unity of the unprincipled, only a false unity opposed to Tradition. Thus, incompatible and incoherent viewpoints can ally with each other on the basis of anti-principle alone. An explicit \#1 is not required because

\begin{quotex}
the locusts have no king, yet but they attack in formation. 
\flright{\textsc{Proverbs 30:27}}
\end{quotex}


Unprincipled themselves, the villagers can only view \#6 as an agent of destruction, hence as something to be opposed. There is no possibility of change within the system.


\flrightit{Posted on 2010-08-30 by Cologero}

